% ============================================================================
% AERIS Paper - Section 8: Conclusion
% Target Journal: Applied Intelligence (APIN) - Springer
% Version: 2026-01-18 (Data Authenticity Corrected)
% ============================================================================

\section{Conclusion}
\label{sec:conclusion}

This paper presents AERIS, a hierarchical routing protocol designed to fill
a critical gap in the WSN protocol design space: \textbf{large-scale
deployments requiring 100\% PDR with robustness under dynamic topologies}.

\subsection{Summary of Contributions}

All contributions below are supported by actual experimental measurements:

\textbf{C1. Scale Reliability}: AERIS maintains \textbf{100\% PDR} at
scales up to 500 nodes, while LEACH degrades to 98.68\% (statistically
significant, $p < 0.001$, Cohen's $d = 1.89$).

\textbf{C2. Energy Efficiency vs LEACH}: AERIS achieves \textbf{18.5\%
energy reduction} compared to LEACH (82.1mJ vs. 100.7mJ at 100 nodes).

\textbf{C3. Honest Trade-off Analysis}: We provide transparent experimental
comparison showing that PEGASIS achieves approximately \textbf{2$\times$
better energy efficiency} than AERIS (41.9mJ vs. 82.1mJ). We explicitly
acknowledge that for energy-critical applications, PEGASIS remains optimal.

\textbf{C4. Robustness}: AERIS maintains 100\% PDR under 30\% node churn
and 40\% regional failure scenarios.

\textbf{C5. Gateway Mechanism Analysis}: Ablation studies reveal the
gateway coordination mechanism as the dominant contributor (Hedges'
$g = 10.09$), while CAS shows negligible independent effect ($g \approx 0$).

\subsection{Honest Positioning}

AERIS is \textbf{not} a universal replacement for existing protocols:

\begin{itemize}
    \item For \textbf{energy-critical} applications (environmental monitoring,
    agricultural sensing), \textbf{PEGASIS} achieves 50\% energy savings
    and should be preferred.

    \item For \textbf{small-scale} deployments ($<$100 nodes), all protocols
    achieve 100\% PDR. Use \textbf{LEACH} for simplicity.

    \item For \textbf{large-scale} applications ($>$200 nodes) requiring
    100\% PDR with energy efficiency better than LEACH, \textbf{AERIS}
    is a viable option.
\end{itemize}

\subsection{Key Findings}

The experimental analysis reveals several important findings:

\begin{enumerate}
    \item The \textbf{Gateway coordination mechanism} is the primary
    contributor to AERIS's reliability (Hedges' $g = 10.09$), while the
    CAS module shows negligible independent effect ($g = 0.00$).

    \item AERIS provides value as a middle-ground between LEACH (simple
    but PDR degrades) and PEGASIS (optimal energy but higher computational
    overhead).

    \item The PDR stability at scale comes from the hierarchical gateway
    structure that provides redundant communication paths.
\end{enumerate}

\subsection{Reproducibility}

All source code, experimental configurations, and raw result data are
released as open source to enable independent verification and facilitate
future research.

\subsection{Future Directions}

Based on our findings, we identify priority directions for future work:

\begin{enumerate}
    \item \textbf{Latency measurement}: Implement hop counting to enable
    actual latency measurement and validate theoretical complexity analysis.

    \item \textbf{Gateway optimization}: Given the Gateway component's
    dominant contribution, future research should focus on adaptive
    gateway mechanisms and energy-aware rotation strategies.

    \item \textbf{Hardware validation}: Deploy AERIS on real TelosB/CC2650
    hardware to validate simulation results and measure actual
    computational overhead.
\end{enumerate}

\subsection{Closing Remarks}

AERIS demonstrates that achieving reliable packet delivery at scale requires
architectural trade-offs. Rather than claiming universal superiority,
we advocate for \textbf{honest, application-specific protocol selection}
that acknowledges each approach's strengths and limitations. We hope
this work contributes not only the AERIS protocol but also a framework
for transparent WSN protocol evaluation that enables practitioners to
make informed deployment decisions.

\textbf{Data Authenticity Statement}: All experimental results reported in
this paper are actual measurements from simulation experiments with 200
independent runs per configuration. Results have been validated using
Welch's t-test with Holm-Bonferroni correction.
